% =====================================================================
%  Pomocná makra
% =====================================================================

% --- poznámky "co sem napsat" ----------------------------------------
% \TODOtrue  = poznámky jsou vidět (žlutě zvýrazněný text)
% \TODOfalse = poznámky zmizí (PŘED ODEVZDÁNÍM)
\newif\ifTODO
\TODOfalse

% zvýraznění textu přes více řádků (soulutf8 = soul s podporou UTF-8)
\usepackage{soulutf8}
\sethlcolor{yellow!45}
% příkazy, které se v \hl{} smí objevit (berou se jako celek)
\soulregister\texttt7
\soulregister\textbf7
\soulregister\emph7
\soulregister\ref7
\soulregister\dots0
\soulregister\soubor7
\soulregister\fce7
\soulregister\tok7

\newcommand{\TODO}[1]{%
  \ifTODO
    \par\noindent{\small\hl{\textbf{TODO:} #1}}\par
    \typeout{TODO: zbyva vyplnit poznamku v dokumentaci}%
  \fi}
\newcommand{\TODOinline}[1]{%
  \ifTODO{\small\hl{#1}}\typeout{TODO: inline}\fi}

% --- titulní strana --------------------------------------------------
% \clen{Jméno Příjmení}{xlogin00}{VUT ID}{procenta}
\newcommand{\clen}[4]{#1 & \texttt{#2} & #3 & #4\,\% \\}
\newcommand{\vedouci}[4]{\textbf{#1} & \texttt{\textbf{#2}} & \textbf{#3} & \textbf{#4\,\%} \\}

% --- identifikátory, soubory -----------------------------------------
\newcommand{\soubor}[1]{\texttt{#1}}     % \soubor{lexer.c}
\newcommand{\fce}[1]{\texttt{#1()}}       % \fce{lexer_next} -> lexer_next()
\newcommand{\tok}[1]{\texttt{#1}}         % \tok{TOKEN\_IDENTIFIER}

% --- gramatika -------------------------------------------------------
\newcommand{\nt}[1]{\ensuremath{\langle}\textit{#1}\ensuremath{\rangle}} % neterminál
\newcommand{\tm}[1]{\texttt{#1}}          % terminál
\newcommand{\eps}{\ensuremath{\varepsilon}}
\newcommand{\pravidlo}{\ensuremath{\rightarrow}}

% pravidlo: \item ... \label{r:x}, v LL tabulce \ref{r:x}
\newlist{gramatika}{enumerate}{1}
\setlist[gramatika]{label=\arabic*.,ref=\arabic*,itemsep=0pt,leftmargin=2.5em}

% --- konečný automat (TikZ) -----------------------------------------
% konvence IFJ: koncový stav = dvojitý kruh, "+" = jeden ze znaků
\usetikzlibrary{automata}
\tikzset{
  ka/.style={
    >={Stealth[round,length=2mm]}, shorten >=1pt, auto,
    font=\footnotesize,
    every state/.style={draw, thick, minimum size=7.5mm, inner sep=0pt,
                        font=\small, fill=white},
    every edge/.style={draw, ->, thin},
    every edge quotes/.style={font=\footnotesize, inner sep=1.5pt},
    initial text={}, initial distance=6mm,
  },
  % souřadnice = pixely původních PNG z diagrams.io
  px/.style={x=0.0098cm, y=-0.0115cm},
  hrana/.style={->, >={Stealth[round,length=2mm]}, shorten >=1pt,
                rounded corners=4pt},
  lbl/.style={font=\footnotesize, inner sep=1.5pt},
  tok/.style={font=\scriptsize\ttfamily, text=black!70},
  pozn/.style={font=\scriptsize\itshape, text=black!60},
}
% symboly v popiscích hran (babel czech dělá z " aktivní znak!)
\newcommand{\kaUv}{\texttt{\textquotedbl}}            % "
\newcommand{\kaBs}{\texttt{\textbackslash}}           % \
\newcommand{\kaPod}{\texttt{\_}}                      % _
\newcommand{\kaPis}{\texttt{a..z}\,+\,\texttt{A..Z}}  % písmena
\newcommand{\kaCis}{\texttt{0..9}}                    % číslice
\newcommand{\kaKrome}[1]{\ensuremath{\Sigma-\{}#1\ensuremath{\}}}

% --- precedenční tabulka -------------------------------------------
\newcommand{\shift}{\ensuremath{<}}
\newcommand{\reduce}{\ensuremath{>}}
\newcommand{\equal}{\ensuremath{=}}
\newcommand{\chyba}{}
